% Copyright 2026 icekylinx and GPT-6 Astra Pro. Apache-2.0.
% Curated from sat_research_checkpoint.zip: paper/integral_full_rectangular_transfer.tex
% Complete full-U mathematical deduction retained with U expanded locally; historical formal-build claims and all weaker numerical/fallback corollaries omitted.
\subsection{The full rectangular curve from integral character profiles}
\label{sec:integral-full-rectangular}

We combine a recent integral coefficient interface with the weighted character-profile analysis of a separate source.  Zhou and
Wang~\cite{zhouwang2026} construct an integer tensor semiring,
prove an integral convolution bound and shared-leg separation,
and derive a universal circuit theorem for the symmetric
rectangular curve.  Their public Lean endpoint covers that
symmetric curve.  The following deduction applies the weighted
profile argument of Amanov, van den Berg, and
Zuiddam~\cite[Sections~4.5--4.6]{amanov2026} to the integral
ingredients, before the optional pairing of characters in
Zhou and Wang's numerical analysis. Their mathematical note also gives a weaker, unformalized bound for
three arbitrary sorted dimensions: $A+B$ when $B\ge2C$, and
$A+B+(2C-B)^2/(4(A+C-B))$ otherwise. The deduction here
strengthens that existing universal three-shape bound to the full
curve $U$; it is not a claim to be the first universal-ring treatment
of arbitrary rectangular shapes.

\begin{theorem}[Additional-source universal integer rectangular curve]
\label{thm:integral-full-rectangular}
Retain the integral spectrum, convolution, separation, and
sector constructions of~\cite{zhouwang2026}.  For every fixed
$A\ge B\ge C\ge0$ and $\epsilon>0$, there is one family of
ordered integer bilinear circuits of size
\[
 O\bigl(n^{U(A,B,C)+\epsilon}\bigr)
\]
for the corresponding rectangular matrix product, valid in
every associative unital ring.  The family and its bound are
chosen before the ring.  The function is explicitly
\[
 U(A,B,C)=\begin{cases}
 A+B,&C(A+B)\le AB,\\[2pt]
 \dfrac{(AB+BC+CA)^2}{4ABC},&C(A+B)>AB,
 \end{cases}
\]
where the first branch includes $C=0$, including $A=B=C=0$.  Integer scalar operations can be
implemented using $0,1,-1$ with the same slackened exponent.
\end{theorem}

\begin{proof}
We spell out the exact interface needed from the integral
source.  Let $\mathcal S_{\mathbb Z}$ contain zero and the
integer three-tensors that restrict to the unit tensor by
integer linear maps, modulo mutual restriction.  Direct sum
and tensor product give its semiring operations.  Its integer
restriction preorder satisfies the Strassen hypotheses, and
the source proves the duality
\[
 \widetilde R_{\mathbb Z}(T)
   =\sup_{\lambda\in\Delta(\mathcal S_{\mathbb Z})}\lambda(T)
 \qquad(T\in\mathcal S_{\mathbb Z}).
\]
All tensors used below have a unit restriction.  In particular,
polynomial convolution, oriented dot products, the sector sums,
and the coefficient-extraction tensors lie in this semiring.
There is no assertion that every nonzero integer tensor lies
in it: a scalar multiple such as $2xyz$ need not do so.

For a character $\lambda$, let $p_i\in[0,1]$ be the exponent
of the oriented dot product whose singleton leg is $i$.
The usual integer reindexings give
$\lambda(B_i(q))=q^{p_i}$ and factor the matrix tensor into
its three oriented dot products.  Write
\[
 C(r,s)=\sum_{0\le i<r,\,0\le j<s}x_i y_j z_{i+j}.
\]
The integral source supplies the following statements for
each individual character, without first averaging it with a
permuted character:
\begin{enumerate}
\item $1\le\lambda(C(r,s))\le r+s-1$, and the same for
every permutation of its legs.  The upper bound comes from
the integral asymptotic-rank convolution lemma, using
Vandermonde adjugates, algebra norms, and a finite B\'ezout
combination after tensor powering.  It is not interpolation
with divisions in the target ring.
\item Character values are monotone under integer coordinate
restrictions and under the displayed integral degenerations.
For the latter, coefficient extraction after tensor powering
has only polynomial rank overhead.
\item If a matched-sector tensor $T$ shares leg $i$ and has
sectors $T_j$, then
\[
 \log\lambda(T)\ge p_i H(q)+\sum_jq_j\log\lambda(T_j)
\]
for each probability vector $q$.  For $p_i>0$, optimization
over $q$ gives
\[
 \lambda(T)^{1/p_i}\ge\sum_j\lambda(T_j)^{1/p_i}.
\]
\item The unimodular determinant filtration degenerates
$C(r,s)\otimes B_1(2)$ into the matched sectors
$C(r,s+1)$ and $C(r,s-1)$, for $s\ge2$.
The factor $B_1(2)$ has value $2^{p_1}$, so the preceding
inequality yields exactly
\[
 \lambda(C(r,s+1))^{1/p_1}
 +\lambda(C(r,s-1))^{1/p_1}
 \le2\lambda(C(r,s))^{1/p_1}.
\]
The factor two is therefore the required discrete-concavity
constant, not a loss that needs to be removed.
\item The integer three-sector construction degenerates
$C(r,3s+r-1)$ into two copies of $C(r,s)$ and one copy with
the last two legs exchanged, all sharing the first leg.
Its middle-branch first-index reversal is an integer
permutation.  Applying the preceding entropy inequality gives
the two coupled tripling inequalities.  Every leg permutation
has the same properties.
\end{enumerate}

These are precisely the character-profile hypotheses used in
Lemma~4.10 and Sections~4.5--4.6 of~\cite{amanov2026}.
To make their field independence explicit, fix a first
coordinate and two corresponding positive row profiles
$g,h$, raised to the power $1/p_i$ when $p_i>0$.
They are discretely concave and satisfy
\[
 2g(s)+h(s)\le g(3s+r-1),\qquad
 2h(s)+g(s)\le h(3s+r-1).
\]
The weighted increment inequality of
\cite[Lemma~4.13]{amanov2026} is a real-variable consequence
of these two facts.  For positive integers $u,v,d,e$ it says
\[
 \left(\frac{g(u+d)}{g(u)}\right)^e
 \left(\frac{h(v+e)}{h(v)}\right)^d
 \ge
 \left(\frac{D}{D-4de/(d+e)}\right)^{d+e},
 \qquad D=u+v+d+e+r-1.
\]
It follows by bounding both increments below by a common
future secant slope and applying weighted AM--GM; no operation
on the coefficient ring occurs.

For positive integers $x_1,x_2,x_3$, put
$S=x_1+x_2+x_3$ and
$Q=p_1x_2x_3+p_2x_1x_3+p_3x_1x_2$.
Apply the weighted inequality to the six permuted convolution
profiles on the integer ray $(nx_1,nx_2,nx_3)$, pairing the
permutations that fix a first leg.  With
$\lambda^\pi(T)=\lambda(\pi\cdot T)$, their weighted product is
\[
 G_n=\prod_{\pi\in S_3}
 \lambda^\pi\!\left(C(nx_{\pi(1)},nx_{\pi(2)})\right)^{x_{\pi(3)}}.
\]
It obeys
\[
 \left(\frac{n+1}{2}\right)^{8Q/S}
       \le G_n\le(nS)^{2S}.
\]
The lower bound is the telescoping product of the six
increment bounds; the upper bound uses only the first
interface item.  When a parameter $p_i$ is zero, monotonicity
supplies the corresponding factor one, exactly as in
\cite[Proposition~4.14]{amanov2026}.
Letting $n$ grow, then using homogeneity and continuity, proves
\begin{equation}\label{eq:integral-character-quadratic}
 4(p_1x_2x_3+p_2x_1x_3+p_3x_1x_2)
       \le(x_1+x_2+x_3)^2
 \quad\text{for all }x_i\ge0.
\end{equation}

The weighted matrix character exponent is
$Ap_2+Bp_3+Cp_1$ under a consistent leg labeling.
If $C(A+B)\le AB$, substitute
$(x_1,x_2,x_3)=(A+B,B,A)$ into
\eqref{eq:integral-character-quadratic}, and use $p_1\ge0$;
the result is $Ap_2+Bp_3+Cp_1\le A+B$.
In the other branch, substitute $(AB,BC,CA)$ to obtain
\[
 Ap_2+Bp_3+Cp_1
       \le\frac{(AB+BC+CA)^2}{4ABC}=U(A,B,C).
\]
Thus every integer character satisfies the full three-shape
bound.  Integral spectral duality supplies the same
asymptotic integer-rank bound.

Finally fix positive slack before selecting one finite integer
rank witness.  For rational shape exponents choose a common
base so that all three dimensions are integer powers, then
take a sufficiently large fixed tensor block.  Recursing on
this ordered bilinear integer scheme costs
$O((t+1)\rho^t)$ at depth $t$: the rank $\rho$ is at least
each of the three input/output coordinate dimensions by
flattening over $\mathbb Q$, and all scalar-linear operations
are charged.  Integer coefficients in this fixed scheme are
implemented by addition and negation.  Padding and rational
approximation from above cover real exponents; continuity of
$U$ and the reserved slack pay for these steps and the
logarithmic recursion factor.  The resulting integer circuit
is chosen before its coefficient ring and preserves the
left-before-right order of each multiplication.
\end{proof}
